\documentclass[12pt]{ximera}
\title{Exam 3}
\begin{document}
\begin{abstract}
Exam 3, covering Chapter 3: matching fair division methods to continuous and discrete
assets, which methods leave someone at exactly a fair share, the method of markers on a
thirteen-item array, lone-divider on a half-veggie, half-beef sub, and sealed bids for
three roommates' chores.
\end{abstract}
\maketitle

\section*{Exam 3}

\begin{problem}
Identify which fair division methods could possibly apply to each situation.

\begin{prompt}
\textbf{(a)} Three siblings need to divide their childhood videogame collection fairly.

\begin{selectAll}
\choice{Lone-Divider Method}
\choice[correct]{Sealed Bids Method}
\choice{Lone-Chooser Method}
\choice[correct]{Method of Markers}
\end{selectAll}

\begin{feedback}
Individual games are indivisible, so the assets are \emph{discrete}: Sealed Bids or
the Method of Markers.
\end{feedback}
\end{prompt}

\begin{prompt}
\textbf{(b)} A $100$-acre property consisting of $15$ acres of beach, $55$ acres of flat
terrain, and $30$ acres of hilly terrain is to be split among four investors.

\begin{selectAll}
\choice[correct]{Lone-Divider Method}
\choice{Sealed Bids Method}
\choice[correct]{Lone-Chooser Method}
\choice{Method of Markers}
\end{selectAll}

\begin{feedback}
Land can be divided into pieces of any size, so the asset is \emph{continuous}:
Lone-Divider or Lone-Chooser. The three terrain types simply make it a mixed
continuous asset, like a half-veggie, half-meat sandwich.
\end{feedback}
\end{prompt}
\end{problem}

\begin{problem}
In which of the following fair division methods is there \emph{always} at least one
player who comes out with exactly their minimum fair share?

\begin{selectAll}
\choice[correct]{Lone-Divider Method}
\choice{Sealed Bids Method}
\choice[correct]{Lone-Chooser Method}
\choice{Method of Markers}
\end{selectAll}

\begin{feedback}
Lone-Divider and Lone-Chooser. In both, the (first) divider cuts pieces they consider
exactly equal, so they end with exactly $\frac1N$ of the value in their own eyes. In
Sealed Bids everyone ends strictly above a fair share once the surplus is split, and in
the Method of Markers the leftover items can lift everyone above it.
\end{feedback}
\end{problem}

\begin{problem}
Players A, B, and C have marked where they would divide the following $13$ items into
three shares of equal value to them.

\begin{center}
\includegraphics[width=\textwidth]{markersArray.png}
\end{center}

The markers fall at A: $4$, $9$; B: $3$, $8$; C: $4$, $7$.

\begin{prompt}
\textbf{(a)} Identify each player's three fair shares. Enter the last item of each
player's first and second shares.

A: items $1$--$\answer{4}$, $5$--$\answer{9}$, and $10$--$13$

\smallskip
B: items $1$--$\answer{3}$, $4$--$\answer{8}$, and $9$--$13$

\smallskip
C: items $1$--$\answer{4}$, $5$--$\answer{7}$, and $8$--$13$

\begin{feedback}
Each marker ends one share, so a player's three shares run from the start to the first
marker, from there to the second marker, and from there to the end.
\end{feedback}
\end{prompt}

\begin{prompt}
\textbf{(b)} Complete the fair division using the method of markers.

Player $\answer{B}$ gets items 1--3, player $\answer{C}$ gets items 4--7, and player
$\answer{A}$ gets items 10--13.

\smallskip
The surplus is items $\answer{8}$ and $\answer{9}$.

\begin{feedback}
First markers: A $4$, B $3$, C $4$; B's is leftmost, so B takes $1$--$3$. Second markers
of A and C: $9$ and $7$; C's is leftmost, so C takes $4$--$7$. A, last, takes everything
after its final marker: $10$--$13$. Items $8$ and $9$ are left over.
\end{feedback}
\end{prompt}
\end{problem}

\begin{problem}
Three friends --- Matt, Kean, and Hope --- want to share an $18$-inch sub that is half
vegetarian (V) and half roast beef (B), using the lone-divider method. The veggie half
is on the left and the beef half on the right. Their preferences:
\begin{itemize}
\item Matt is vegan, so doesn't value the beef half at all;
\item Kean likes both halves equally well;
\item Hope considers the veggie half worth $60\%$ of the total and the beef worth $40\%$.
\end{itemize}

\begin{prompt}
\textbf{(a)} Write each friend's value system, as fractions.

Matt --- V: $\answer{1}$, B: $\answer{0}$ \quad
Kean --- V: $\answer{1/2}$, B: $\answer{1/2}$ \quad
Hope --- V: $\answer{3/5}$, B: $\answer{2/5}$

\begin{feedback}
Matt puts all of his value on the veggie half; Kean splits evenly; Hope's $60\%$ and
$40\%$ are $\frac35$ and $\frac25$.
\end{feedback}
\end{prompt}

\begin{prompt}
\textbf{(b)} If Hope is the divider, what three shares does she cut? Enter fractions.

$s_1=\answer{5/9}$ of V \qquad
$s_2=\frac49$ of V $+$ $\answer{1/6}$ of B \qquad
$s_3=\answer{5/6}$ of B

\begin{hint}
Each share must be worth $\frac13$ to Hope. The veggie half alone is worth $\frac35$ to
her --- more than $\frac13$ --- so the first cut falls inside the veggie.
\end{hint}

\begin{feedback}
The first share is a fraction $x$ of V with $\frac35x=\frac13$, so $x=\frac59$. The
remaining $\frac49$ of V is worth $\frac49\cdot\frac35=\frac4{15}$, so the second share
needs $\frac13-\frac4{15}=\frac1{15}$ more, from beef: $\frac25y=\frac1{15}$ gives
$y=\frac16$. The third share is the other $\frac56$ of B, worth
$\frac56\cdot\frac25=\frac13$. \checkmark

On the $18$-inch sub, the cuts fall $5$ inches from the veggie end and $1.5$ inches into
the beef.
\end{feedback}
\end{prompt}

\begin{prompt}
\textbf{(c)} What is each share worth to each friend? Enter fractions.

Matt: $s_1=\answer{5/9}$, $s_2=\answer{4/9}$, $s_3=\answer{0}$

\smallskip
Kean: $s_1=\answer{5/18}$, $s_2=\answer{11/36}$, $s_3=\answer{5/12}$

\smallskip
Hope: $s_1=\answer{1/3}$, $s_2=\answer{1/3}$, $s_3=\answer{1/3}$

\begin{feedback}
Matt values only the veggie: $\frac59$, $\frac49$, $0$.

Kean values by size: $\frac59\cdot\frac12=\frac5{18}$;
$\frac49\cdot\frac12+\frac16\cdot\frac12=\frac29+\frac1{12}=\frac{11}{36}$;
$\frac56\cdot\frac12=\frac5{12}$.

Hope, the divider, values each at exactly $\frac13$.

Fair shares ($\ge\frac13$): Matt $\set{s_1,s_2}$; Kean only $\set{s_3}$; Hope any.
\end{feedback}
\end{prompt}

\begin{prompt}
\textbf{(d)} Can the shares be fairly assigned without cutting further?

\begin{multipleChoice}
\choice[correct]{Yes}
\choice{No}
\end{multipleChoice}

\begin{feedback}
Yes. Kean needs $s_3$, and nobody else needs it; Matt can take either $s_1$ or $s_2$;
Hope takes whatever remains.
\end{feedback}
\end{prompt}

\begin{prompt}
\textbf{(e)} Describe one fair division of the sub using these shares.

\begin{freeResponse}
\end{freeResponse}

\begin{feedback}
For example: Kean takes $s_3$ ($\frac5{12}$ to him), Matt takes $s_1$ ($\frac59$), and
Hope takes $s_2$ ($\frac13$). Swapping Matt and Hope --- Matt $s_2$ ($\frac49$), Hope
$s_1$ --- is the other fair option.
\end{feedback}
\end{prompt}
\end{problem}

\begin{problem}
Roommates Wei, Yan, and Ali are using the method of sealed bids to split weekly chores.
They've estimated the cost of each chore to themselves:

\begin{center}
\begin{tabular}{|r|r|r|r|}\hline
 & \textbf{Wei} & \textbf{Yan} & \textbf{Ali}\\\hline
Tidying & $-\$8$ & $-\$4$ & $-\$10$\\
Dishes & $-\$12$ & $-\$14$ & $-\$15$\\
Laundry & $-\$11$ & $-\$9$ & $-\$6$\\
Groceries & $-\$70$ & $-\$80$ & $-\$60$\\
Trash & $-\$4$ & $-\$10$ & $-\$2$\\\hline
\end{tabular}
\end{center}

\begin{prompt}
\textbf{(a)} Who is assigned each chore?

Tidying: $\answer{Yan}$ \quad Dishes: $\answer{Wei}$ \quad Laundry: $\answer{Ali}$ \quad
Groceries: $\answer{Ali}$ \quad Trash: $\answer{Ali}$

\begin{feedback}
Each chore goes to whoever minds it least (the least negative bid): Tidying Yan ($-4$),
Dishes Wei ($-12$), Laundry Ali ($-6$), Groceries Ali ($-60$), Trash Ali ($-2$).
\end{feedback}
\end{prompt}

\begin{prompt}
\textbf{(b)} Fill in the rest of the table.

Total value: Wei $\answer{-105}$, Yan $\answer{-117}$, Ali $\answer{-93}$

\smallskip
Fair share: Wei $\answer{-35}$, Yan $\answer{-39}$, Ali $\answer{-31}$

\smallskip
Allocation value: Wei $\answer{-12}$, Yan $\answer{-4}$, Ali $\answer{-68}$

\smallskip
Initial settlement (positive = pays in): Wei $\answer{23}$, Yan $\answer{35}$,
Ali $\answer{-37}$

\smallskip
Total surplus: $\$\answer{21}$, so each roommate's share is $\$\answer{7}$

\begin{feedback}
Fair shares are a third of each total. Wei's dishes ($-12$) are $\$23$ better than his
fair share of $-35$, so he pays $\$23$; Yan's tidying ($-4$) beats $-39$ by $\$35$; Ali's
three chores ($-68$) are $\$37$ worse than $-31$, so Ali receives $\$37$. The system
holds $23+35-37=\$21$, or $\$7$ per roommate.
\end{feedback}
\end{prompt}

\begin{prompt}
\textbf{(c)} Altogether, how much does each roommate owe or receive, in addition to
their chores? Enter a positive number for money paid, negative for money received.

Wei: $\answer{16}$ \quad Yan: $\answer{28}$ \quad Ali: $\answer{-44}$

\begin{feedback}
Subtract each \$7 share of the surplus: Wei pays $23-7=\$16$, Yan pays $35-7=\$28$, and
Ali receives $37+7=\$44$. The payments balance: $16+28=44$. \checkmark

Ali takes on the groceries, the most burdensome chore, and is paid well for it; every
roommate ends \$7 better off than a fair share by their own estimates.
\end{feedback}
\end{prompt}
\end{problem}

\end{document}
